\subsection{The complete central parameter of the ramified marked-cusp family}

\begin{statement}{E.70}{Proved}{Finite orbit transitions and full product coordinates for the marked central parameter}
Retain the actual groups and Frobenius of E.67:
$B=\Pi/I_Q$, $D=\mathbb Z_3[[B]]$,
$C=\widehat{\bigwedge}^2D$,
$\theta=F_Q|_B$, and $q=|k_v|$.
Put $s=v_3(q-1)\ge1$ and
\[
          \mathcal Z_Q=C/(q^2\theta_2-1)C,\qquad
          \theta_2=\widehat{\bigwedge}^2\theta_* .
 \tag{E70.1}
\]
There is a homeomorphism of compact $\mathbb Z_3$-modules
\[
                    \mathcal Z_Q
                         \simeq\prod_{n\ge1}\mathbb Z_3 .
 \tag{E70.2}
\]
The actual noncanonical coordinates and their inverse
are constructed in (E70.10)--(E70.13) from finite orbit
characters and coherent division, retaining the
Frobenius and every transition map.

More precisely, for a characteristic open subgroup
$U\subset B$, let $\Omega_U$ be the unordered pairs
of distinct elements of $B/U$. For each $\theta$-orbit
$O\subset\Omega_U$, choose an oriented representative
$(g_O,h_O)$. Its orbit length $L_O$ is a power of three.
Transport this orientation by $\theta$; it returns
with the same orientation. Then the actual finite
quotient is
\[
 \begin{split}
 Z_U&=\operatorname{coker}\!
   \left(q^2\theta_2-1:
       \bigwedge^2\mathbb Z_3[B/U]
          \longrightarrow\bigwedge^2\mathbb Z_3[B/U]\right)\\
    &\simeq\bigoplus_{O\in\Omega_U/\theta}
           \mathbb Z_3/(q^{2L_O}-1)
     \simeq\bigoplus_O\mathbb Z/3^{\,s+v_3(L_O)},\\
 \mathcal Z_Q&\xrightarrow{\ \sim\ }\varprojlim_U Z_U .
 \end{split}
 \tag{E70.3}
\]
For $U'\subset U$, an oriented finer orbit generator
maps to zero when its endpoints coincide modulo $U$.
Otherwise its image is
\[
             e_{O'}\longmapsto
                    \epsilon q^{-2j}e_O ,
 \qquad \epsilon\in\{1,-1\},
 \tag{E70.4}
\]
where its reduced ordered pair is $\epsilon$ times
the $j$th transported orientation of the parent orbit.
This is a homomorphism between the displayed finite
cyclic factors; all finite maps in (E70.3) are
surjective. The same formula applies to every
$\theta$-stable open normal subgroup; its finite
quotient is dominated by a characteristic one.

Consequently the ramified marked-cusp inertia locus
of E.67 has the complete coordinates
\[
 (\mathbb Z_3\setminus\{0\})\times
       \mathbb Z_3/(q-1)\mathbb Z_3\times
                       \prod_{n\ge1}\mathbb Z_3 .
 \tag{E70.5}
\]
The actual full image is precisely the zero sequence
in the last factor, with the representatives and
unique full lifts of E67.5--E67.7.
The word E67.9 is detected by its specified finite
moment character; its finite torsion quotient is
compatible with the torsion-free compact module
(E70.2) through that actual quotient map.

At $q=4$, the actual rank-nine quotient
$B/B'B^3$ gives
\[
      Z_{B'B^3}\simeq
        (\mathbb Z/3)^{2390391}
                 \oplus(\mathbb Z/9)^{63770004}.
 \tag{E70.6}
\]
There is also an explicit primitive arithmetic
quotient $B\twoheadrightarrow\mathbb Z_3$ on which
$\theta$ is multiplication by four. Its positional
level $3^n$ gives the fully evaluated finite system
\[
      Z^{\rm line}_n\simeq
         \bigoplus_{e=0}^{n-1}
             (\mathbb Z/3^{e+1})^{\,4\cdot3^e-1},
 \tag{E70.7}
\]
with the literal signed phase transitions (E70.4).
The program evaluates all of these transitions
through $n=6$, as well as coherent character lifts.
These are actual quotients of (E70.3).

\emph{Manuscript: \S51.16, folios733v--736v,
scans172--178, formulas(371)--(379), Part E(c),(d);
the marked subgroup and its actual quotient maps
are exactly those in E.41 and E.67.}
\end{statement}

\begin{proof}
\emph{The exact finite orbit system.}
By E.67 the normal-kernel centre is the compact
module
$C=\varprojlim_U\bigwedge^2\mathbb Z_3[B/U]$.
The closure of $\theta$ is pro-three. Therefore its
permutation of $B/U$ has three-power order, and every
orbit in $\Omega_U$ has odd length.
An ordered pair cannot return with its endpoints
exchanged: that would give a nontrivial permutation
of order two induced by a group of odd order.
Thus, on each orbit, the transported vectors
\[
       e_j=\delta_{\theta^jg_O}\wedge
                     \delta_{\theta^jh_O}
                   \quad(0\le j<L_O)
\]
give an integral basis with
$\theta_2e_j=e_{j+1}$, indices read modulo $L_O$.
The relation $q^2e_{j+1}-e_j$ gives
$[e_j]=q^{-2j}[e_0]$, and the final relation is
$(q^{2L_O}-1)[e_0]=0$.
Eliminating the intermediate basis vectors uses
unit integral coefficients. This proves the first
two lines of (E70.3), including its literal generator.
The lifting-the-exponent identity gives
$v_3(q^{2L_O}-1)=s+v_3(L_O)$.

The finite operators $q^2\theta_2-1$ are injective:
a power of the permutation is the identity, while
$q^{2m}-1$ is a nonzero integer for $m>0$.
Their short exact cokernel sequences are compact.
The same finite-intersection lifting argument as
in E66.10 makes their inverse limit exact.
It yields exactly the final map in (E70.3).
Its denominator is the compact image of the
continuous operator on $C$.

Under reduction of finite coset sets a wedge
becomes zero when its two endpoints coalesce.
Otherwise its orientation and position in the
parent orbit give (E70.4). The parent orbit
length divides the finer length, so its cyclic
relation annihilates the image of the finer
relation. Each parent wedge lifts to a finer
wedge; the resulting column is a unit times its
parent generator. Hence each displayed finite
cokernel map is surjective.

\emph{No three-torsion survives in the actual compact quotient.}
The exterior module has the exact alternating
realization
\[
 C\xrightarrow{\ \sim\ }
   \{\mu\in\mathbb Z_3[[B\times B]]:
                         \operatorname{swap}_*\mu=-\mu\},
 \quad
 \delta_g\wedge\delta_h\longmapsto
           \delta_{(g,h)}-\delta_{(h,g)}.
 \tag{E70.8}
\]
The inverse is one half of the natural exterior
projection. The idempotent
$(1-\operatorname{swap})/2$ is integral.
At each finite positional level its image is a
direct summand of a finite free lattice.
Compact exact limits of the multiplication-by-three
sequences show that reducing (E70.8) modulo three
gives exactly the alternating direct summand of
$\mathbb F_3[[B\times B]]$.

E67.13 proves $B^{\theta^m}=1$ for every $m>0$,
using the actual finite-field weights and the
separated integral lower-central filtration.
The only periodic point of $\theta\times\theta$
on $B\times B$ is therefore $(1,1)$.
An invariant $\mathbb F_3$-measure on this compact
space is supported at that point. Indeed, above
any coarse orbit whose inverse image avoids
$(1,1)$, compactness forces all refining orbit
lengths to grow beyond every fixed three-power
multiple. Trace compatibility sums equal
coefficients around each refined cycle, and
eventually makes the coefficient of the coarse
orbit a multiple of three. It is therefore zero.
This proves the asserted support statement on
every clopen set.

Since $q^2\equiv1\pmod3$, an element of the kernel
of $q^2\theta_2-1$ on $C/3C$ would be such an
invariant alternating measure. The Dirac measure
at $(1,1)$ is symmetric, and two is invertible,
so that kernel is zero.
Suppose $3x=(q^2\theta_2-1)y$ in $C$.
Reduction modulo three gives $y=3y_1$; the
torsion-freeness of $C$ then gives
$x=(q^2\theta_2-1)y_1$.
We have proved
\[
                         \mathcal Z_Q[3]=0.
 \tag{E70.9}
\]
The same argument applies to all powers of three.
Thus the compact pro-three module is torsion-free.

\emph{Countably infinitely many independent directions.}
The finitely generated pro-three group $B$ has
a countable cofinal system of characteristic open
subgroups. The positional and coefficient
presentations make $C$ and $\mathcal Z_Q$
metrizable. Their continuous character groups
are countable unions of finite groups.

Choose a strictly decreasing cofinal sequence
$U_1\supset U_2\supset\cdots$ of characteristic
open subgroups, and set $A_n=U_n\setminus U_{n+1}$.
There are continuous alternating characters
with values in $\mathbb F_3$ defined by
\[
 \lambda_n(\delta_g\wedge\delta_h)=
       1_{A_n}(g)1_{U_{n+1}}(h)
           -1_{A_n}(h)1_{U_{n+1}}(g).
 \tag{E70.10}
\]
They are $\theta_2$-invariant and hence annihilate
$q^2\theta_2-1$ modulo three.
For $g_m\in A_m$, their values on
$\delta_{g_m}\wedge\delta_1$ are exactly
$\delta_{nm}$.
Thus the order-three continuous character
space of $\mathcal Z_Q$ has countably infinite
dimension. This argument uses actual
Frobenius-stable clopen sets, including their
centred limit at the identity.

\emph{Coherent finite characters and the actual product map.}
Write $\mathcal A=\varinjlim_U Z_U^\vee$, where
$Z_U^\vee=\operatorname{Hom}(Z_U,\mathbb Q_3/\mathbb Z_3)$.
Every transition in this direct system is the
dual injection of (E70.4).
The order-three characters of each finite cyclic
orbit factor enumerate $\mathcal A[3]$.
Choose a basis $\alpha_1,\alpha_2,\ldots$ of that
space by successive finite linear algebra in this
enumeration. To compare any finite list of
characters, pull them to one common finite
positional level; injectivity makes that a
literal finite computation.

Every finite character has an explicit division
by three at a sufficiently fine level.
For the finitely many parent orbits on which it
is nonzero, choose a finer characteristic
subgroup so that each surviving finer orbit
above them has length divisible by three times
the parent length. Such a subgroup exists
uniformly: the inverse image of a parent
unordered distinct pair is compact, contains
no diagonal point, and contains no periodic
point of $\theta\times\theta$.
The orbit-length argument used above applies.
Put $t_O=s+v_3(L_O)$.
For a character with value $a_O/3^{t_O}$ on
the parent generator, define its divided
character on a finer generator by
\[
       e_{O'}\longmapsto
          \frac{\epsilon q^{-2j}a_O}{3^{t_O+1}}
                  \pmod{\mathbb Z_3};
 \tag{E70.11}
\]
set the value to zero on generators mapping
to zero or to a zero parent coordinate.
The finer exponent is at least $t_O+1$, so
these are legitimate characters of the actual
finite cyclic factors. Multiplying them by
three gives the exact pullback of the original
character. The unit in the numerator is the
same phase unit as in (E70.4), now evaluated
to the required higher precision.

Starting with $\alpha_{i,1}=\alpha_i$, this
procedure chooses coherent characters
$\alpha_{i,n}$ of order $3^n$ with
$3\alpha_{i,n+1}=\alpha_{i,n}$.
For $z\in\mathcal Z_Q$, define
\[
 \Phi(z)_i=(a_{i,n})_{n\ge1},\qquad
 \alpha_{i,n}(z)=a_{i,n}/3^n\pmod{\mathbb Z_3},
 \quad a_{i,n}\in\mathbb Z/3^n.
 \tag{E70.12}
\]
Here $a_{i,n}$ is the unique residue with
the displayed character value.
Coherence makes this a continuous
$\mathbb Z_3$-linear map to
$\prod_i\mathbb Z_3$.

The character chains form a direct sum of
Pruefer three-groups and span $\mathcal A$.
For independence, take a finite relation at
common order $3^n$ and multiply by $3^{n-1}$.
Independence of the $\alpha_i$ makes every
coefficient divisible by three. Repeating
proves that every coefficient is zero at
the original order.
For spanning, take a character of order
$3^n$, express its $3^{n-1}$ multiple in
the basis $\{\alpha_i\}$, and subtract the
corresponding finite sum of $\alpha_{i,n}$.
The remainder has smaller order; induction
finishes the expression.

This also gives an explicit inverse to
(E70.12). For a prescribed sequence
$(u_i)\in\prod_i\mathbb Z_3$, evaluate
$\alpha_{i,n}$ as $u_i/3^n$ modulo
$\mathbb Z_3$. The just-proved direct sum
makes this a well-defined homomorphism on
$\mathcal A$. Restrict it to the character
group of each $Z_U$ and use the elementary
duality of its displayed finite cyclic
factors. The resulting $z_U\in Z_U$ are
compatible with (E70.4), so give
\[
              (u_i)\longmapsto
                   (z_U)_U\in\varprojlim_U Z_U.
 \tag{E70.13}
\]
All finite evaluations use only the character
expressions obtained by the order induction
and the finite linear algebra above.
Equations (E70.12)--(E70.13) are inverse
continuous maps of compact modules.
This proves (E70.2) with actual product
coordinates and their inverse.

\emph{Arithmetic finite examples and all scalar transitions.}
For $q=4$, E67.22 gives the actual matrix $F_B$.
Modulo three, $F_B-I$ has rank two and square
zero. Thus its permutation of $\mathbb F_3^9$
has order three and exactly $3^7$ fixed points.
An unordered pair is fixed exactly when both
points are fixed, because a permutation of
odd order cannot exchange just those two
points. Therefore the numbers of pair orbits
of lengths one and three are
\[
 \binom{3^7}{2}=2390391,\qquad
  \frac{\binom{3^9}{2}-\binom{3^7}{2}}3=63770004 .
 \tag{E70.14}
\]
The finite presentation (E70.3) proves
(E70.6), without enumerating a much larger
matrix than these orbit counts require.

In the E.37 basis, the primitive functional
\[
       f_2^*+e_{A_0}^*+e_{A_1}^*-e_{B_2}^*
           :B^{\rm ab}\longrightarrow\mathbb Z_3
 \tag{E70.15}
\]
annihilates the filled cusp vector and
satisfies $\ell F_B=4\ell$.
It takes value one on $e_{A_0}$, so is
surjective. Composing with $B\to B^{\rm ab}$
gives the actual arithmetic quotient to
$\mathbb Z_3$ claimed in the statement.
Pushforward of measures and exterior
powers gives corresponding surjections
from $C$ and $\mathcal Z_Q$ to the scalar
orbit systems.

On $\mathbb Z/3^n$, multiplication by four
has period $3^{k-1}$ on a pair whose
minimum valuation is $n-k$.
The number of unordered distinct pairs of
that valuation is
$4\cdot3^{2k-2}-3^{k-1}$: subtract the
diagonal from the two-coordinate valuation
stratum, then divide by two.
Consequently there are exactly
$4\cdot3^{k-1}-1$ orbits of that length.
This proves (E70.7) for every $n$.
Each surviving finer orbit has three times
the parent length; those whose endpoints
coalesce map to zero. Hence the character
division (E70.11) can be carried out at
the next level in this scalar system.

The exact program \texttt{a\_central\_orbits.py}
checks the actual primitive functional and
the rank-nine mod-three matrix, every
oriented scalar pair through level729,
all finite coboundaries and transitions,
the finite order-three basis extensions,
378 coherent character divisions, and
the independent annular characters.
Its retained output distinguishes the full
rank-nine finite quotient from the explicit
scalar arithmetic quotient throughout.
No uncomputed higher Frobenius word is
assigned a numerical value.
\end{proof}
